% docs/slides/preamble.tex -- 五份 deck 共享
\documentclass[aspectratio=169,11pt]{beamer}

% ==== 中文与字体 ====
\usepackage[UTF8,fontset=none]{ctex}
\usepackage{fontspec}

\IfFontExistsTF{PingFang SC}{%
  \setCJKmainfont{PingFang SC}[BoldFont=PingFang SC,ItalicFont=PingFang SC,AutoFakeSlant=0.15]
  \setCJKsansfont{PingFang SC}[ItalicFont=PingFang SC,AutoFakeSlant=0.15]
  \setCJKmonofont{PingFang SC}
}{%
  \IfFontExistsTF{Source Han Sans SC}{%
    \setCJKmainfont{Source Han Sans SC}
    \setCJKsansfont{Source Han Sans SC}
  }{%
    \IfFontExistsTF{Noto Sans CJK SC}{%
      \setCJKmainfont{Noto Sans CJK SC}
      \setCJKsansfont{Noto Sans CJK SC}
    }{%
      \setCJKmainfont{Heiti SC}
      \setCJKsansfont{Heiti SC}
    }
  }
}

\IfFontExistsTF{Helvetica Neue}{\setmainfont{Helvetica Neue}}{%
  \IfFontExistsTF{Inter}{\setmainfont{Inter}}{}}
\IfFontExistsTF{JetBrains Mono}{\setmonofont{JetBrains Mono}[Scale=0.85]}%
  {\setmonofont{Menlo}[Scale=0.85]}

% ==== 颜色 ====
\usepackage{xcolor}
\definecolor{cMain}{HTML}{1F4E79}
\definecolor{cAccent}{HTML}{E07A1F}
\definecolor{cGray}{HTML}{595959}
\definecolor{cMute}{HTML}{F4F4F4}
\definecolor{cOK}{HTML}{2E7D32}
\definecolor{cBad}{HTML}{B23A3A}
\definecolor{cWarn}{HTML}{C8A300}

% ==== 主题 ====
\usepackage{tcolorbox}
\usetheme{metropolis}
\metroset{progressbar=frametitle,numbering=fraction,block=fill}
\setbeamercolor{frametitle}{bg=cMain,fg=white}
\setbeamercolor{progress bar}{fg=cAccent,bg=cMute}
\setbeamercolor{alerted text}{fg=cAccent}
\setbeamercolor{title}{fg=cMain}

% metropolis 默认尝试 Fira Sans，未安装时 fallback 到 Latin Modern Sans。
% 这里在主题加载后显式覆盖，强制使用 macOS 上一定存在的 Helvetica Neue，
% 跟中文 PingFang 风格一致；找不到则按顺序尝试 Inter / Arial。
\IfFontExistsTF{Helvetica Neue}{%
  \setsansfont{Helvetica Neue}%
}{%
  \IfFontExistsTF{Inter}{\setsansfont{Inter}}{%
    \IfFontExistsTF{Arial}{\setsansfont{Arial}}{}}%
}

% ==== TikZ ====
\usepackage{tikz}
\usetikzlibrary{positioning,arrows.meta,fit,shapes.geometric,calc,%
  decorations.pathmorphing,backgrounds}

\tikzset{
  node-box/.style={draw=cMain,thick,rounded corners=2pt,
    minimum height=8mm,minimum width=18mm,align=center,font=\small},
  node-state/.style={draw=cMain,thick,circle,minimum size=14mm,
    align=center,font=\small},
  node-ok/.style={draw=cOK,thick,fill=cOK!10},
  node-warn/.style={draw=cWarn,thick,fill=cWarn!15},
  node-bad/.style={draw=cBad,thick,fill=cBad!10},
  % 色盲友好的状态机三态：蓝 (正常/Closed) / 橙 (告警/Open) / 灰 (中间/HalfOpen)
  % 避免 red-green 二元区分，改用 hue + lightness 双通道。
  node-state-normal/.style={draw=cMain,thick,fill=cMain!12},
  node-state-alert/.style={draw=cAccent,thick,fill=cAccent!18},
  node-state-transition/.style={draw=cGray,thick,fill=cGray!18},
  arrow/.style={-{Stealth[length=2mm]},thick,draw=cMain},
  arrow-async/.style={-{Stealth[length=2mm]},thick,dashed,draw=cGray},
  arrow-bad/.style={-{Stealth[length=2mm]},thick,draw=cBad},
  arrow-alert/.style={-{Stealth[length=2mm]},thick,draw=cAccent},
  lifeline/.style={dashed,thick,draw=cGray},
}

% ==== 代码 ====
\usepackage{listings}
\lstdefinelanguage{Go}{
  morekeywords={break,case,chan,const,continue,default,defer,else,
    fallthrough,for,func,go,goto,if,import,interface,map,package,range,
    return,select,struct,switch,type,var,nil,true,false,iota},
  morekeywords=[2]{string,int,int64,uint,uint64,bool,byte,error,context},
  sensitive=true,
  morecomment=[l]//,morecomment=[s]{/*}{*/},
  morestring=[b]",morestring=[b]`,
}
\lstdefinelanguage{LuaRedis}{
  morekeywords={and,break,do,else,elseif,end,false,for,function,if,in,
    local,nil,not,or,repeat,return,then,true,until,while},
  morekeywords=[2]{redis,call,KEYS,ARGV,tonumber,HGET,HSET,GET,SET,
    DECR,DECRBY,INCR,INCRBY,EXPIRE,PEXPIRE,ZADD,ZCARD,ZREMRANGEBYSCORE,
    EVAL,EVALSHA},
  sensitive=true,
  morecomment=[l]{--},
  morestring=[b]",morestring=[b]',
}

\lstdefinestyle{base}{
  basicstyle=\ttfamily\scriptsize,
  numbers=left,numberstyle=\tiny\color{cGray},
  keywordstyle=\color{cMain}\bfseries,
  keywordstyle=[2]\color{cAccent},
  stringstyle=\color{cOK},
  commentstyle=\color{cGray}\itshape,
  backgroundcolor=\color{cMute},
  frame=single,framerule=0pt,
  showstringspaces=false,
  breaklines=true,
  tabsize=2,
  columns=fullflexible,
}
\lstset{style=base}

% ==== 三线表与附录页码 ====
\usepackage{booktabs}
\usepackage{appendixnumberbeamer}

% ==== 页脚来源标注 ====
\newcommand{\srcnote}[1]{{\color{cGray}\scriptsize\itshape #1}}

% 关键论点框
\newcommand{\keypoint}[1]{%
  \begin{tcolorbox}[colback=cMute,colframe=cMain,boxrule=0.5pt,arc=2pt]%
  #1\end{tcolorbox}}


\title{搜索是流量入口}
\subtitle{从 LIKE 搜不到"iPhone 13" 到 ES + outbox 增量索引}
\author{gomall · 业务侧 deck}
\date{}

\begin{document}

\maketitle

\begin{frame}{今日大纲}
\tableofcontents
\end{frame}

% ============================================================
\section{搜不到 = GMV 直接掉}
% ============================================================

\begin{frame}{从首页到下单的真实漏斗}
\begin{tikzpicture}[node distance=4mm and 8mm]
  \node[node-box] (home) {首页\\100\%};
  \node[node-box,right=of home] (sch)  {搜索\\$\sim$60\%};
  \node[node-box,right=of sch,node-warn]  (lst)  {结果列表\\$\sim$45\%};
  \node[node-box,right=of lst]  (det)  {详情页\\$\sim$18\%};
  \node[node-box,right=of det,node-ok]  (ord)  {下单\\$\sim$3\%};
  \draw[arrow] (home) -- (sch);
  \draw[arrow] (sch)  -- (lst);
  \draw[arrow] (lst)  -- (det);
  \draw[arrow] (det)  -- (ord);
  \draw[arrow-bad,bend left=25] (sch.north) to node[above,font=\tiny]{搜不到 $\to$ 直接流失} (ord.north);
\end{tikzpicture}
\vspace{2mm}
\begin{itemize}
  \item 进店用户里大约 60\% 会用搜索框。
  \item "搜索 $\to$ 结果列表"这一跳是漏斗最陡的一段。
  \item 这一跳 miss，用户既不去详情页，也不去类目页，直接流失。
\end{itemize}
\srcnote{routes/router.go:40 /product/search 是这条漏斗的入口}
\end{frame}

\begin{frame}{"搜不到"在业务上等于什么}
\begin{itemize}
  \item \textbf{C 端}：用户搜一个词没结果，多数不会再换词，直接退出 App。
  \item \textbf{商家}：上架了商品但搜不到，等同于"商品在仓库里没人知道"。
  \item \textbf{运营}：投放素材引来流量，关键词在站内搜不到，等于把广告费白烧。
  \item \textbf{客服}：每天都会接到"明明有为什么搜不到"的工单。
\end{itemize}
\vspace{2mm}
\keypoint{搜索不是"功能"，是\textbf{流量入口}。入口 miss 一次 = 一单变零单。}
\end{frame}

% ============================================================
\section{LIKE 为什么不够}
% ============================================================

\begin{frame}{反面教材：用户搜"苹果手机"}
\begin{tikzpicture}[node distance=8mm and 14mm]
  \node[node-box] (q) {用户输入\\\textit{苹果手机}};
  \node[node-box,right=of q,node-bad] (sql) {\texttt{name LIKE}\\\texttt{'\%苹果手机\%'}};
  \node[node-box,right=of sql] (row) {商品标题\\\textit{iPhone 13 Pro}};
  \node[node-box,below=of sql,node-bad] (miss) {0 命中};
  \draw[arrow] (q)   -- (sql);
  \draw[arrow] (sql) -- (row);
  \draw[arrow-bad] (sql) -- (miss);
\end{tikzpicture}
\vspace{2mm}
\begin{itemize}
  \item 用户脑子里是中文，商家上架习惯写英文型号。
  \item LIKE 不知道"苹果 == Apple == iPhone"，子串匹配必然 miss。
  \item 这不是用户输入错，是\textbf{搜索系统的责任}。
\end{itemize}
\srcnote{repository/db/dao/product.go:80--98 当前 LIKE 实现}
\end{frame}

\begin{frame}{LIKE 的另一个问题：跑不动}
\begin{tabular}{@{}lr@{}}
\toprule
指标 & 数值 \\
\midrule
product 表行数         & 956K 行 \\
product 表磁盘大小     & 364 MB \\
LIKE \texttt{'\%X\%'} 索引可用？ & 否（前缀通配） \\
全表扫预估延迟         & 秒级 \\
\bottomrule
\end{tabular}
\vspace{3mm}
\begin{itemize}
  \item \texttt{LIKE '\%kw\%'} 前缀是通配符，B+ 树索引用不上，必扫全表。
  \item 96 万行 + 360 MB 数据，单次搜索就吃满一颗 CPU。
  \item 对照：\texttt{/product/show} 走主键，压测 \textbf{62K RPS / p95 3ms}（同样的库）。
\end{itemize}
\srcnote{stressTest/REPORT.md \S 数据规模 + 链路压测表}
\end{frame}

\begin{frame}{LIKE 缺的三件事}
\begin{tabular}{@{}lll@{}}
\toprule
能力      & LIKE 有？ & 业务后果 \\
\midrule
分词       & 无 & "苹果手机" 搜不到 "iPhone" \\
同义词     & 无 & "拖鞋" 搜不到 "凉拖" \\
相关性排序 & 无 & 名称命中 vs 描述命中 等权排 \\
\bottomrule
\end{tabular}
\vspace{3mm}
\begin{itemize}
  \item LIKE 把搜索退化成"子串包含"，跟用户的搜索意图基本错位。
  \item 即便加索引、加全文索引（MySQL FULLTEXT），中文分词支持依然差。
  \item 这是\textbf{业务上换 ES 的根本原因}，不是为了"用一个高大上的组件"。
\end{itemize}
\end{frame}

% ============================================================
\section{ES 给业务的三件事}
% ============================================================

\begin{frame}{ES 替 LIKE 做什么}
\begin{tikzpicture}[node distance=6mm and 12mm]
  \node[node-box]                       (q)   {用户输入\\\textit{苹果手机}};
  \node[node-box,right=of q,node-warn]  (an)  {分词\\\textit{苹果} / \textit{手机}};
  \node[node-box,right=of an]           (syn) {同义词\\Apple / iPhone};
  \node[node-box,right=of syn,node-ok]  (sc)  {打分\\name${}^3$ title${}^2$ info};
  \node[node-box,below=of sc]           (out) {结果列表\\相关性降序};
  \draw[arrow] (q)   -- (an);
  \draw[arrow] (an)  -- (syn);
  \draw[arrow] (syn) -- (sc);
  \draw[arrow] (sc)  -- (out);
\end{tikzpicture}
\vspace{2mm}
\begin{itemize}
  \item \textbf{分词}：中文按词切，不是按字符。
  \item \textbf{同义词}：词典 / 模型把"苹果手机"映射到 Apple / iPhone。
  \item \textbf{相关性}：命中标题加权高于命中描述。gomall 走的是 \texttt{multi\_match} + 字段加权。
\end{itemize}
\srcnote{repository/es/product\_index.go:164--173 multi\_match 查询 / 字段权重}
\end{frame}

\begin{frame}{字段加权背后的业务判断}
\begin{tabular}{@{}llr@{}}
\toprule
字段 & 业务含义 & 加权 \\
\midrule
\texttt{name}  & 商品主名（用户脑中预期匹配） & $\times 3$ \\
\texttt{title} & 商品标题（含型号 / 规格）   & $\times 2$ \\
\texttt{info}  & 商品详情（长文，噪声多）    & $\times 1$ \\
\bottomrule
\end{tabular}
\vspace{3mm}
\begin{itemize}
  \item 用户搜 "iPhone"，结果里"标题就叫 iPhone 13" 应当排在 "info 里出现 iPhone" 前面。
  \item 加权是\textbf{运营调出来的}，不是开发拍的。促销期 \texttt{title} 权重可以临时上调。
  \item 不加权的结果是"卖标签的店铺把 info 塞满热词"反而排第一。
\end{itemize}
\srcnote{repository/es/product\_index.go:170 fields=\{name\^{}3, title\^{}2, info\}}
\end{frame}

\begin{frame}{精度 vs 召回：业务侧怎么挑}
\begin{itemize}
  \item \textbf{召回}：用户搜的，相关商品至少出现一个 $\to$ "搜得到"。
  \item \textbf{精度}：返回的前 10 条里，真正相关的占比 $\to$ "搜得准"。
  \item 这两件事互相挤压。
\end{itemize}
\vspace{2mm}
\begin{tabular}{@{}lll@{}}
\toprule
场景    & 优先 & 原因 \\
\midrule
长尾搜索 (\textit{雕花茶台}) & 召回 & 没结果就直接流失 \\
品牌搜索 (\textit{iPhone}) & 精度 & 太多杂牌山寨页拉低体验 \\
营销关键词 (\textit{618 折扣}) & 精度 & 第一屏决定下单率 \\
\bottomrule
\end{tabular}
\vspace{2mm}
\small gomall 现在统一走 \texttt{multi\_match}，没区分场景。后续要做 query 分类。
\end{frame}

% ============================================================
\section{上架到可搜的 SLA}
% ============================================================

\begin{frame}{商家最常问：我上架的商品多久能搜到}
\begin{itemize}
  \item 这是一个\textbf{业务承诺}，不是技术细节。
  \item 商家上架后立刻去搜索框验证，搜不到就会发工单。
  \item 业务侧承诺：\textbf{上架 5 分钟内可搜}。这个数字写进商家后台 FAQ。
  \item 工程侧要把这个数字落到具体链路里。
\end{itemize}
\vspace{3mm}
\keypoint{业务承诺 $\to$ SLO $\to$ 监控告警阈值。三层是一条线。}
\end{frame}

\begin{frame}{从上架到可搜的完整链路}
\begin{tikzpicture}[node distance=5mm and 10mm]
  \node[node-box]                      (api) {商家上架\\API};
  \node[node-box,right=of api]         (db)  {MySQL\\product 表};
  \node[node-box,right=of db,node-warn](ob)  {outbox\\表};
  \node[node-box,right=of ob]          (mq)  {RabbitMQ\\product.changed};
  \node[node-box,right=of mq]          (cs)  {indexer\\消费者};
  \node[node-box,right=of cs,node-ok]  (es)  {ES\\product 索引};
  \draw[arrow] (api) -- (db);
  \draw[arrow] (db)  -- node[above,font=\tiny]{同事务} (ob);
  \draw[arrow-async] (ob) -- node[above,font=\tiny]{异步} (mq);
  \draw[arrow] (mq)  -- (cs);
  \draw[arrow] (cs)  -- (es);
\end{tikzpicture}
\vspace{2mm}
\begin{itemize}
  \item 商家 API 返回 200，DB 写入完成，但用户还\textbf{搜不到}。
  \item 真正可搜要等 outbox $\to$ MQ $\to$ 消费者 $\to$ ES 这一串走完。
  \item 正常情况整条链路 $<$ 1s。慢时可达数分钟。
\end{itemize}
\srcnote{service/product.go:257--265 outbox 写事件；service/search/indexer.go:32--47 消费者}
\end{frame}

\begin{frame}{延迟分解：5 分钟预算花在哪}
\begin{tabular}{@{}lrr@{}}
\toprule
环节                     & 正常 & 异常 \\
\midrule
DB 写入 + outbox 同事务   & < 50ms & 几百 ms（主从慢） \\
outbox publisher 轮询    & < 1s   & 几秒（积压时） \\
RMQ 投递                 & < 100ms& 几秒（broker 慢） \\
indexer 消费 + 调 ES     & < 200ms& 几秒到分钟（ES 慢） \\
ES refresh\_interval     & 1s     & 1s（默认配置） \\
\midrule
总计                     & $\sim$1s & 5 分钟内基本能覆盖 \\
\bottomrule
\end{tabular}
\vspace{2mm}
\small 业务承诺的 5 分钟，是把最坏情况乘 2 留余量得出的，不是技术上最佳。\\
ES 写入用 \texttt{refresh=false}，refresh\_interval 决定可搜延迟下界。
\srcnote{repository/es/product\_index.go:119 Refresh: "false"}
\end{frame}

\begin{frame}{为什么不直接同步写 ES}
\begin{itemize}
  \item 朴素做法：商家 API 里 \texttt{tx.Commit()} 之后直接 \texttt{es.Upsert}。
  \item 问题一：ES 调用失败，DB 已经提交，状态\textbf{不一致}。
  \item 问题二：ES 调用慢，商家 API 的 p99 被 ES 拖累。
  \item 问题三：ES 挂了，商家完全无法上架。
  \item gomall 用 outbox 把"业务写"和"搜索同步"解耦：
  \begin{itemize}
    \item 商家 API 只依赖 DB，ES 故障不影响上架。
    \item ES 恢复后 outbox 自动追回积压事件。
  \end{itemize}
\end{itemize}
\srcnote{对应技术 deck 04 Outbox：业务 + outbox 同事务，publisher 异步投 MQ}
\end{frame}

\begin{frame}{全量回填：上线 / 灾后兜底}
\begin{itemize}
  \item 增量索引解决的是"\textbf{以后}的商品"，存量商品要靠全量回填。
  \item 触发场景：
  \begin{itemize}
    \item 第一次切 ES：把 956K 行 product 灌进去。
    \item ES 集群重建 / 索引 mapping 改字段。
    \item outbox 长时间积压后人工对账。
  \end{itemize}
  \item gomall 暴露 admin 接口手动触发，按主键游标分批 200。
  \item 不放在客户端入口，避免被误触发把 ES 打挂。
\end{itemize}
\srcnote{routes/router.go:130 admin/search/backfill；service/search/backfill.go:13--43 实现}
\end{frame}

% ============================================================
\section{ES 挂了用户看到什么}
% ============================================================

\begin{frame}{静默降级，而不是 5xx}
\begin{tikzpicture}[node distance=8mm and 12mm]
  \node[node-state,node-state-normal] (n) {正常\\ES 在};
  \node[node-state,node-state-transition,right=of n] (d) {降级\\走 DB};
  \node[node-state,node-state-alert,right=of d] (a) {告警\\运营介入};
  \draw[arrow]       (n) -- node[above,font=\tiny]{ES 超时} (d);
  \draw[arrow-alert] (d) -- node[above,font=\tiny]{持续 > 5min} (a);
  \draw[arrow,bend left=30] (d.south) to node[below,font=\tiny]{ES 恢复} (n.south);
\end{tikzpicture}
\vspace{2mm}
\begin{itemize}
  \item ES 客户端启动失败 $\to$ \texttt{tryInitES} recover $\to$ 进程照常起来。
  \item \texttt{ProductSearch} 检测 \texttt{EsClient == nil} 或 ES 查询报错 $\to$ 走 DB LIKE。
  \item 用户体验：搜得到的词依然搜得到，召回退化但不会 5xx。
\end{itemize}
\srcnote{cmd/main.go:61--69 tryInitES recover；service/product.go:294--323 ES 失败 fall back DB}
\end{frame}

\begin{frame}{降级时用户实际感受是什么}
\begin{tabular}{@{}lll@{}}
\toprule
搜索词        & 正常 (ES) & 降级 (LIKE) \\
\midrule
\textit{iPhone 13}    & 命中 & 命中（标题里有） \\
\textit{苹果手机}      & 命中 & \textbf{0 命中} \\
\textit{白色短袖}      & 命中 (相关) & 命中（顺序乱） \\
\textit{618}           & 命中加权 & 命中 + 慢 \\
\bottomrule
\end{tabular}
\vspace{3mm}
\begin{itemize}
  \item 降级期间长尾词 / 语义词\textbf{真的搜不到}，召回率显著下降。
  \item 但比起返回 5xx，至少漏斗还在转，主流量词不掉。
  \item 业务可以接受这个降级，前提是\textbf{监控能在 5min 内告警}。
\end{itemize}
\end{frame}

\begin{frame}{ES 长时间挂掉的处理流程}
\begin{enumerate}
  \item \textbf{T+0}：ES 告警，搜索接口自动走 DB。
  \item \textbf{T+5min}：值班 SRE 拉群，运营同步"搜索能力降级"。
  \item \textbf{T+15min}：如果是热点词召回严重下降，运营在 banner 提示"搜索能力恢复中"。
  \item \textbf{T+恢复}：ES 起来，outbox 自动追积压事件，无需人工干预。
  \item \textbf{T+恢复后}：监控 \texttt{outbox\_lag} 回到 0，关闭事件。
\end{enumerate}
\vspace{2mm}
\small 整个过程\textbf{不需要}DB 兜底数据搬到 ES，因为 outbox 已经记了。
\srcnote{对应技术 deck 04 Outbox：publisher 失败重试 + at-least-once 投递}
\end{frame}

% ============================================================
\section{遗留风险与多角色视角}
% ============================================================

\begin{frame}{遗留风险：重试可能"老盖新"}
\begin{tikzpicture}[node distance=6mm and 10mm]
  \node[node-box]                (u1) {update v1};
  \node[node-box,right=of u1]    (u2) {update v2};
  \node[node-box,right=of u2]    (mq) {RMQ};
  \node[node-box,right=of mq,node-bad]    (cs) {indexer\\重试};
  \node[node-box,right=of cs]    (es) {ES doc};
  \draw[arrow] (u1) -- (u2);
  \draw[arrow] (u2) -- (mq);
  \draw[arrow] (mq) -- (cs);
  \draw[arrow-bad] (cs) -- node[above,font=\tiny]{读到 v1 又写回} (es);
\end{tikzpicture}
\vspace{2mm}
\begin{itemize}
  \item indexer 拿到 \texttt{product\_id}，再回 DB 读最新行写 ES。
  \item 但 MQ 是 at-least-once，重试时 DB 可能已经被另一条事件改回 v2 $\to$ 实际无害。
  \item 真正风险：两条事件\textbf{并发消费}，老的事件后到 ES，覆盖了新值。
  \item 修法：ES 写入带 \texttt{version} / \texttt{updated\_at}，老版本拒写。\textbf{当前未做，列为已知缺口}。
\end{itemize}
\srcnote{service/search/indexer.go:51--60 handleProductChanged 当前实现}
\end{frame}

\begin{frame}{四个角色看搜索}
\begin{tabular}{@{}llp{5.5cm}@{}}
\toprule
角色   & 关心        & 看哪个数字 \\
\midrule
产品   & 转化率      & 搜索 $\to$ 详情页跳转率 \\
运营   & 关键词收益  & 热搜词 GMV / 0 结果词 top N \\
商家   & 上架延迟    & 上架 $\to$ 可搜 p99 < 5min \\
SRE   & 集群稳定性  & ES 5xx 率 / outbox 积压量 \\
\bottomrule
\end{tabular}
\vspace{3mm}
\small "0 结果词 top N"是运营每天必看的报表：用户搜了但没结果的词，要么补品类要么扩同义词。
\end{frame}

\begin{frame}{回到业务承诺}
\begin{itemize}
  \item 搜到 vs 搜不到：差的不是"算法精度"，是\textbf{有没有人能下到那一单}。
  \item LIKE $\to$ ES 这一步换的不是技术栈，是\textbf{搜索系统对召回的承诺}。
  \item Outbox + 增量索引这套链路换的不是异步框架，是\textbf{对商家"5 分钟可搜"的承诺}。
  \item 静默降级换的不是高可用配置，是\textbf{"ES 挂了用户不掉单"的承诺}。
\end{itemize}
\vspace{2mm}
\keypoint{每一项技术选择，背后都对应一个写在 FAQ 里的业务承诺。}
\end{frame}

% ============================================================
\appendix
% ============================================================

\begin{frame}{业务侧 Q\&A（上）}
\textbf{Q1}：为什么不直接用 MySQL FULLTEXT，要单独引入 ES？\\
\hspace{4mm}\small A：MySQL FULLTEXT 中文分词需要外挂插件，相关性打分能力弱，
而且和 OLTP 在同一个实例上抢资源。ES 独立部署、独立扩缩容，搜索流量再大也不会拖 DB。

\vspace{2mm}
\textbf{Q2}：商家上架后多久能搜到？\\
\hspace{4mm}\small A：业务承诺 5 分钟内。技术链路正常情况 $<$ 1s，承诺数字留 5--10 倍余量给 ES 慢 / RMQ 积压等异常。

\vspace{2mm}
\textbf{Q3}：ES 慢 5 秒不可用怎么办？\\
\hspace{4mm}\small A：读路径自动降级回 DB LIKE，写路径不受影响。
outbox 事件会积压等 ES 恢复后追，最长几分钟，不丢事件。
\end{frame}

\begin{frame}{业务侧 Q\&A（下）}
\textbf{Q4}：降级到 DB 之后用户感受是什么？\\
\hspace{4mm}\small A：精确匹配词（\textit{iPhone 13}）依然搜得到；同义词 / 长尾词（\textit{苹果手机}）会 0 命中。
监控 5 分钟内告警，运营在前台 banner 提示"搜索能力恢复中"。

\vspace{2mm}
\textbf{Q5}：商家改了一个商品名，搜索结果什么时候更新？\\
\hspace{4mm}\small A：和上架同链路，正常 $<$ 1s。改名 + outbox + indexer + ES。

\vspace{2mm}
\textbf{Q6}：已知缺口里那个"老盖新"什么时候修？\\
\hspace{4mm}\small A：单商品并发改名的概率低，目前没出现客诉。修法是在 ES 写入加 \texttt{version}
字段，规划在下一个版本随相关性调优一起做。
\end{frame}

\begin{frame}{代码位置一览}
\begin{description}
  \item[搜索路由入口]\srcnote{routes/router.go:40}
  \item[全量回填 admin 入口]\srcnote{routes/router.go:130}
  \item[DB 路径 LIKE 实现]\srcnote{repository/db/dao/product.go:80--98}
  \item[ES 多字段查询 + 加权]\srcnote{repository/es/product\_index.go:156--209}
  \item[ProductSearch 降级逻辑]\srcnote{service/product.go:294--325}
  \item[outbox 写事件]\srcnote{service/product.go:257--265}
  \item[indexer 消费者]\srcnote{service/search/indexer.go:17--60}
  \item[全量回填实现]\srcnote{service/search/backfill.go:13--43}
  \item[ES 启动静默降级]\srcnote{cmd/main.go:61--69 + initialize/search.go:13--31}
\end{description}
\vspace{2mm}
\keypoint{配套技术 deck：04-outbox-saga.tex（outbox 与最终一致性）}
\end{frame}

\end{document}
